% !TeX root = ../main.tex

\section*{Supplementary Methods}

\subsection*{Supplementary Methods 1. Agent construction, role prompts and meeting protocol}

The workflow was implemented as a sequence of structured human--agent meetings
rather than as a single autonomous model call. Agent objects were instantiated
from a common role schema containing a title, expertise statement, project goal,
stage-specific role description and model identifier. The project configuration
recorded \texttt{gpt-5.4-2026-03-05} as the model identifier for the
role-specialized agents, with two discussion rounds used for team meetings and
three independent iterations used for stages that required merged summaries or
aggregation. The recorded discussion files, merged summaries and plotting
artifacts were retained under phase-specific folders in the project workspace.

The core workflow contained a Principal Investigator agent, three domain-scientist
agents and a Scientific Critic. The Principal Investigator led meetings,
sequenced agendas, synthesized role-specific inputs and generated the retained
stage-level recommendations. The Maillard Reaction Chemistry and Process
Scientist evaluated glycation chemistry, donor reactivity, dry-state processing,
temperature--time trade-offs, water-activity effects, ultrasound or pretreatment
burden and AGE/byproduct risk. The Allergenicity and Immunoassay Interpretation
Scientist assessed whether lower antibody-binding signals were likely to reflect
biologically meaningful immunoreactivity reduction or instead assay-format,
aggregation, extractability or epitope-presentation effects. The Experimental
Design and Evidence Synthesis Scientist mapped tested and untested condition
space, preserved comparator logic, identified high-information comparisons and
converted open optimization questions into compact wet-lab matrices. The
Scientific Critic challenged overextension, missing controls, unsupported
mechanistic claims and decisions that relied on reduction magnitude without
appropriate comparator support. A separate Research Synthesis and Optimization
role was used for literature-facing synthesis and evidence-structuring tasks,
including the structured extraction workflow described in Supplementary
Methods~2.

The team-selection stage was run before the main optimization meetings. The
human researcher supplied the project objective, the need to optimize
$\beta$-LG glycation under practical constraints, and the requirement that the
team include complementary expertise in glycation chemistry, immunological
interpretation and experimental design. The Principal Investigator then selected
the three domain-scientist roles listed above and excluded itself from the
domain-scientist team, preserving the intended distinction between meeting lead
and specialist contributors.

Two meeting modes were used. In team meetings, the agenda, previous merged
summaries and any relevant data packets were provided to the Principal
Investigator and the active agent team. The Principal Investigator first framed
the decision problem, then each team member responded in a fixed order. The
Principal Investigator synthesized the first-round inputs, posed follow-up
questions or requested clarification, and a second round of responses was
collected. At the end of the meeting, the Principal Investigator produced a
detailed summary, a stage-level decision or recommendation and any requested
answer to the agenda. In individual or aggregation meetings, one agent, usually
the Principal Investigator, received several independent discussion summaries and
merged them into a single retained output. This pattern was used when the stage
required a consensus artifact, such as role selection, project specification or
tool-screening synthesis.

Information transfer between stages was explicit. Each new phase received the
current project background, the stage agenda, relevant prior merged summaries and
the structured artifacts needed for that decision, such as literature-extraction
records, sequence-prior outputs, run-sheet fields or returned wet-lab results.
The output of each phase was stored as \texttt{discussion\_1.md},
\texttt{discussion\_2.md}, \texttt{discussion\_3.md} where applicable, together
with a \texttt{merged.md} and/or \texttt{merged.json} file. These merged outputs
were treated as the auditable state passed to downstream phases, rather than as
private model memory.

Human oversight was present at four decision points. First, the human researcher
defined the protein system, endpoint, feasible processing envelope and practical
constraints. Second, the human researcher reviewed whether proposed tools were
applicable to non-enzymatic glycation rather than enzymatic glycosylation and
whether they were feasible to execute for this project. Third, the human
researcher executed or coordinated wet-lab work and returned structured result
packets instead of allowing the agents to infer unavailable experimental values.
Fourth, the human researcher confirmed ambiguous metadata when source workbooks
omitted fields or used shorthand labels. The agents therefore performed
intermediate synthesis, criticism and planning, but they did not replace wet-lab
execution, feasibility judgment or final human acceptability checks.

\subsection*{Supplementary Methods 2. Literature corpus assembly, evidence extraction and harmonization}

The literature-derived evidence layer comprised 23 full-text studies available in
the project corpus and successfully processed using the structured extraction
workflow. The corpus covered non-enzymatic glycation or Maillard-type modification
of $\beta$-lactoglobulin ($\beta$-LG) and related food proteins,
processing-assisted glycation, and immunoreactivity- or allergenicity-related
measurements. It was assembled as a task-specific evidence corpus rather than as a
systematic review or exhaustive meta-analysis. The included studies provided
identifiable experimental runs, principal reaction or processing conditions,
associated comparators and at least one relevant outcome measurement.

The 23 studies were processed using a fixed Evidence Curator instruction
set. The agent was instructed to act as a structured information extractor and to
return JSON Lines (JSONL) only, without headings, explanations or Markdown. The
recording unit was one paper, one experimental run and one measured metric per row.
Consequently, a run evaluated using more than one immunoreactivity-related metric
generated multiple records sharing the same run identifier. Each record contained
exactly 23 fields: \texttt{paper\_id}, \texttt{run\_id},
\texttt{comparator\_run\_id}, \texttt{protein}, \texttt{sugar},
\texttt{protein\_sugar\_ratio}, \texttt{mode}, \texttt{temperature\_C},
\texttt{time\_h}, \texttt{pH\_or\_aw}, \texttt{pretreatment\_used},
\texttt{pretreatment\_method}, \texttt{pretreatment\_params},
\texttt{metric\_name}, \texttt{assay\_name}, \texttt{unit},
\texttt{direction}, \texttt{treated\_value}, \texttt{control\_value},
\texttt{reduction\_pct\_reported}, \texttt{reduction\_pct\_basis},
\texttt{reduction\_pct\_source} and \texttt{reduction\_pct\_note}. The field
definitions and controlled vocabulary are provided in Supplementary Table~1, and
the complete machine-readable records are available in the project data repository.

Identifiers and condition fields were extracted conservatively. The
\texttt{paper\_id} used the first author and publication year where available.
The \texttt{run\_id} preserved the group or sample name used in the source article,
and \texttt{comparator\_run\_id} was required to refer to a run from the same paper.
Reaction mode was restricted to \texttt{wet} or \texttt{dry}. Solution-state
conditions were recorded using an explicit pH label, whereas dry-state conditions
were recorded using an explicit relative-humidity or water-activity label to avoid
ambiguous expressions. Pretreatment was separated from the principal glycation
reaction; its presence was stored as a Boolean value, its method as a text label,
and reported settings as key--value pairs. Information absent from the source was
stored as \texttt{null} and was not guessed or reconstructed.

For every metric-level record, \texttt{treated\_value} denoted the value associated
with the current \texttt{run\_id}, whereas \texttt{control\_value} denoted the value
of its matched comparator. Comparator records themselves had a null comparator
identifier and contained only their own measured value. Metric direction was encoded
as either \texttt{higher\_is\_more\_allergenic} or
\texttt{higher\_is\_less\_allergenic}. Binding, optical-density, activation,
release and related signal measurements were generally assigned the former
direction, whereas inhibitory-concentration measurements such as IC50 were assigned
the latter when a larger value indicated weaker immune binding. Where the direction
was not explicit in the article, the inference was documented in the note field.

The initial extraction prioritized reduction percentages explicitly reported by
the source article. Values were accepted from unambiguous statements such as
``reduced by $X$\%'', ``reduced to $Y$\% of control'' or ``$Z$\% remained''; the
latter two forms were converted to $100-Y$ or $100-Z$, respectively. A fold change
was converted only when the source clearly defined its reduction-oriented meaning.
Every accepted value retained its basis metric, a concise source locator comprising
the PDF page and relevant Results, figure or table location, and a note describing
any conversion. If only a percentage was reported, absolute treated and control
values remained null. This extraction produced 68 treatment/comparison records with
directly reported reduction percentages.

To increase comparator-relative coverage without conflating reported and derived
values, a separate harmonization step calculated an additional reduction-oriented
effect only when both treated and matched control values were numerically available.
For metrics in which a higher value indicated greater immunoreactivity, the
harmonized reduction was calculated as
\[
R = \left(1-\frac{x_{\mathrm{treated}}}{x_{\mathrm{control}}}\right)\times100.
\]
For metrics such as IC50 in which a higher value indicated lower immunoreactivity,
the direction-aware reduction was calculated as
\[
R = \left(1-\frac{x_{\mathrm{control}}}{x_{\mathrm{treated}}}\right)\times100.
\]
Negative values were retained because they indicated a change in the direction
opposite to reduced immunoreactivity. This step added 86 calculable records, yielding
154 treatment/comparison records with either directly reported or direction-aware
reduction-oriented effect information. The remaining records were not assigned a
numeric reduction when either member of the treated--comparator pair was absent.

Before downstream use, the JSONL records were checked for the presence of a
\texttt{paper\_id}, exact field count, valid enumeration values, non-ambiguous
pH/relative-humidity/water-activity notation, within-paper comparator links,
consistent treated and control assignments, and concise source-location strings.
The validated evidence layer contained 350 metric-level records representing 207
unique experimental runs: 280 treatment/comparison records and 70 comparator/control
records. Because the studies differed in protein system, reaction backbone,
assay format and comparator definition, these records were not interpreted as a
single directly rankable dataset. Comparator category, metric direction and local
experimental context were preserved during all subsequent evidence synthesis.

\subsection*{Supplementary Methods 3. Computational tool screening and sequence-prior generation}

The computational-tool stage was designed as a feasibility and prioritization
layer for non-enzymatic $\beta$-LG glycation, not as a complete simulation of
Maillard chemistry or as a quantitative model of IgG-binding reduction. The
initial agent discussion considered a broad structural triage stack, including
experimental $\beta$-LG A structures or conformational ensembles, solvent
accessibility analysis, local protonation or reactivity annotation, focused
small-sugar docking and epitope-prediction tools. The intended abstraction was
``glycation opportunity'' near epitope-relevant regions: whether lysine-centred
surface patches were plausible modification sites and whether those sites
overlapped with sequence regions likely to contribute to immune recognition.

Candidate tools were then screened by the human researcher before execution. The
screening first separated tools applicable to non-enzymatic glycation from tools
or workflows designed primarily for enzymatic glycosylation, glycan-chain
annotation or mature glycoform modelling. Tools in the latter category were not
advanced because the present study did not seek to predict enzymatically installed
glycan structures. A second feasibility screen removed analyses that would require
validated pretreatment-specific conformational ensembles, dry-state donor
diffusion models or calibrated covalent-reaction parameters. In particular,
docking and molecular-dynamics workflows were retained only as future heuristic
options, because ordinary non-covalent docking does not model Schiff-base
formation, Amadori rearrangement, dry-state diffusion or heterogeneous multi-site
Maillard reaction kinetics. The final executed computational layer was therefore
restricted to two reproducible sequence-level priors: BepiPred-3.0 for B-cell
epitope propensity and NetGlycate for lysine-level non-enzymatic glycation
susceptibility. The full screening logic is summarized in Supplementary Table~2.

Both sequence-prior analyses used the bovine $\beta$-LG variant A sequence from
PDB 1BSQ chain A. The input sequence contained 162 residues:
\begin{quote}
\footnotesize
\texttt{LIVTQTMKGLDIQKVAGTWYSLAMAASDISLLDAQSAPLRVYVEELKPTPEGD}\\
\texttt{LEILLQKWENGECAQKKIIAEKTKIPAVFKIDALNENKVLVLDTDYKKYLLFC}\\
\texttt{MENSAEPEQSLACQCLVRTPEVDDEALEKFDKALKALPMHIRLSFNPTQLEEQCHI}
\end{quote}
The archived FASTA input and residue-level output files are included in
the project data repository.

BepiPred-3.0 was run on the 1BSQ chain A sequence to obtain residue-level
B-cell epitope propensity scores. Because the exported file indexed residues by
sequence order rather than by a separate position column, residue positions were
assigned according to the input FASTA sequence. The BepiPred-3.0 score ranged
from 0.0417 to 0.2929 across the 162 residues, and the threshold displayed with
the archived interactive output was 0.1512. For visualization and interpretation,
three high-priority sequence windows were annotated: residues 61--69, 110--115
and 130--148. These windows were selected because they contained local maxima in
the BepiPred-3.0 output and overlapped regions that were useful for discussing
epitope-proximal glycation opportunity. The lysine residues K69 and K77 were
labelled in the corresponding figure because they were among the highest-scoring
lysine positions in the BepiPred-3.0 output, with scores of 0.2910 and 0.2903,
respectively.

NetGlycate was run on the same 162-residue sequence to obtain lysine-level
non-enzymatic glycation susceptibility scores. In the archived output, lysine
positions with scores greater than or equal to 0 were treated as sequence-level
glycation-prone annotations for plotting. This identified K47, K70, K14, K100
and K101 as positive-score lysines, with NetGlycate scores of 0.925, 0.818,
0.804, 0.803 and 0.801, respectively. These values were interpreted only as
relative sequence-level susceptibility flags. They were not used to infer
site-specific glycation yields, glycation stoichiometry or expected IgG-binding
reduction.

The sequence-prior outputs were used qualitatively during project specification.
BepiPred-3.0 provided a map of epitope-propensity regions, and NetGlycate
identified lysines with sequence-level glycation plausibility. The co-occurrence
of epitope-propensity windows and nearby lysine annotations, such as the
61--69/K69--K70 region, supported attention to epitope-proximal glycation but
did not establish causal site masking. No aggregate computational optimization
score was calculated, no donor or reaction condition was selected solely from
the sequence-prior outputs, and no structural docking result was used as evidence
for the final wet-lab recommendation. All final condition choices required
literature evidence, returned wet-lab data and human feasibility adjudication.

\subsection*{Supplementary Methods 4. Project specification, question-structured run-sheet generation and availability audit}

Project specification was carried out after the structured evidence layer and the
sequence-level prior layer had been generated. The purpose of this stage was to
convert the broad $\beta$-LG glycation condition space into a limited set of
experimentally answerable questions, rather than to produce an exhaustive factorial
screen. Meeting inputs included the 23-study structured evidence layer, the
BepiPred-3.0 and NetGlycate sequence-prior outputs, previous wet-lab benchmark
information, feasibility constraints supplied by the human researcher and the
preceding agent-meeting summaries. The optimization target was defined as reduced
IgG-binding immunoreactivity under a dry-state $\beta$-LG variant A system, with
the practical aim of identifying a short, interpretable and experimentally
feasible high-reduction window.

The planning discussion organized the decision space around five comparison axes:
donor identity, the choice between 55~$^\circ$C and 60~$^\circ$C, the value of
extending reaction time, restrained pentose advancement and the separation of
donor-associated glycation from the common processing backbone. The primary
matrix fixed reaction mode, water activity and protein:sugar ratio to preserve
comparability across donor and temperature comparisons. The fixed sugar-containing
backbone was dry-state glycation at $a_w$ 0.79 with a protein:sugar ratio of
1:2 (w/w). Broad wet-mode screening, pH optimization, water-activity sweeps,
protein:sugar-ratio optimization and alternative pretreatment families were
deferred because adding these axes in the same round would have made the returned
comparisons less interpretable.

The question-structured planning output contained 24 condition rows. The matrix
was not designed as a balanced full factorial screen. Instead, it contained
targeted blocks that encoded specific scientific questions. R01--R08 formed a
matched hexose bridge, comparing lactose, glucose, galactose and mannose at
55~$^\circ$C and 60~$^\circ$C for 4~h under ultrasound-assisted dry-state
processing. R09--R12 were no-ultrasound sentinels for the same four hexose donors
at 55~$^\circ$C for 4~h, intended to test the local contribution of ultrasound
within this matched backbone. R13--R16 provided lactose and glucose time anchors
at 60~$^\circ$C for 2~h and 55~$^\circ$C for 3~h. R17--R20 formed a restrained
pentose branch with arabinose and ribose at 55~$^\circ$C for 2--3~h. R21--R24
were heated no-sugar process controls designed to distinguish donor-associated
glycation from the common heating background.

Comparator protection was specified before result analysis. Native untreated
$\beta$-LG was used as the assay reference for calculating IgG-binding reduction,
but was not treated as a process-matrix run. Heated no-sugar entries were retained
as process-only controls for the corresponding dry-state backbone. Time-anchor
runs were used to judge whether a longer reaction earned its added processing
burden, and no-ultrasound sentinels were used only as condition-matched tests of
the local ultrasound contribution at 55~$^\circ$C for 4~h. These comparator roles
were encoded in the run-sheet rather than inferred after observing the outcomes.

Before formal analysis, each planned condition was matched to the returned wet-lab
records using a condition key consisting of donor, temperature, reaction time,
ultrasound status, reaction mode, water activity, protein:sugar ratio and
no-sugar-control status. The supplemental workbook \textit{ELISA-2026.6.15.xlsx}
was incorporated to complete the previously unavailable no-ultrasound hexose
sentinels, the ribose 55~$^\circ$C, 2~h condition and the 55~$^\circ$C no-sugar
2~h and 3~h controls. When a source workbook did not print one condition field
but the condition identity could be resolved from the sample label and run-sheet
context, the run was retained and marked as a direct match with harmonized
metadata. Specifically, the lactose, glucose, galactose and mannose 4~h rows in
this workbook were treated as no-ultrasound 55~$^\circ$C conditions for R09--R12,
and the \textit{N-U-3h} entry was treated as the 3~h no-sugar ultrasound control
for R22.

All 24 planned conditions had directly attributable results for the primary
IgG-binding readout after incorporation of the supplemental workbook. Therefore,
the availability audit no longer required reindexing from the original planning
identifiers, and R01--R24 were retained as the analysis identifiers. The audit was
applied before effect-size interpretation and did not use reduction magnitude,
donor identity or whether a result was favourable as an inclusion criterion.

The completed run matrix was organized into five functional blocks: a matched 4~h
hexose bridge, no-ultrasound hexose sentinels, benchmark and practical time
anchors, a restrained pentose branch and heated no-sugar process controls. Each
row stored the run identifier, decision block, donor, temperature, time,
ultrasound status, reaction mode, water activity, protein:sugar ratio, control
status, result status, reduction value, source workbook, source label, figure
panel and audit note. The complete 24-condition plan and availability audit are
provided in Supplementary Table~3, and the run-to-source mapping is provided in
Supplementary Table~4, with machine-readable files available in the project data repository.

\subsection*{Supplementary Methods 5. Wet-lab feedback, adjudication and recommendation locking}

Wet-lab feedback was returned to the agent workflow as structured decision input
rather than as a single ranked list of reduction values. Each feedback packet
retained the condition fields used during planning, the primary IgG-binding
reduction value relative to the untreated $\beta$-LG baseline, the available
comparator context, assay-quality notes and any process observations that could
affect interpretation. The Principal Investigator agent received these feedback records together with the preceding project-specification summary, sequence-level prior outputs and human-defined feasibility constraints. Domain agents and the Scientific Critic then evaluated the returned results against the comparison logic that had been specified before execution.

Adjudication followed a comparator-aware hierarchy. Donor-containing runs were
first interpreted relative to the untreated baseline to identify whether a branch
produced a meaningful reduction in IgG binding. Sugar-containing runs were then
considered against same-backbone heated no-sugar process controls where such
controls were available, so that signal loss caused by heating or recovery alone
was not overinterpreted as donor-associated glycation. Temperature comparisons
were evaluated within matched donor and time settings, and time comparisons were
used only when the donor, reaction state, water activity, protein:sugar ratio and
pretreatment backbone were otherwise aligned. The no-ultrasound sentinel rows
R09--R12 enabled a local ultrasound comparison against their matched 55~$^\circ$C,
4~h hexose counterparts, but this comparison was not used to claim a general
ultrasound effect across other donors, temperatures or reaction times. Runs with
incomplete comparator support could inform branch prioritization, but they were
not used to claim a general process effect.

The same marginal-benefit rule used during project specification was retained
during feedback interpretation. For two consecutive time points $t_1$ and $t_2$,
the marginal gain was defined as
\[
G(t_1,t_2)=\frac{R(t_2)-R(t_1)}{t_2-t_1},
\]
where $R(t)$ denotes IgG-binding reduction at reaction time $t$. The quantity was
reported in percentage points per hour (pp h$^{-1}$), representing the absolute
increase in reduction per additional hour rather than a relative percentage
change. Additional time or temperature was considered useful only when the
expected gain was meaningful relative to the added processing burden and
interpretability cost.

The glucose time-course block was used only for this marginal-benefit calibration.
Under the fixed dry-state, ultrasound-assisted framework, glucose reached
70.99\%, 78.75\%, 82.50\%, 84.33\%, 87.91\% and 90.81\% IgG-binding reduction at
3, 6, 9, 12, 18 and 24~h, respectively. These values showed that longer reaction
time could increase reduction, but also that the increase was non-linear and
should not be treated as a uniformly strengthening axis. The glucose curve was
therefore not used as a prediction curve for ribose. Its role was to calibrate
how the agent team interpreted diminishing gains, time burden and the plausibility
of reaching high reduction within a shorter practical window.

Final donor selection was made using an evidence hierarchy rather than a
post-hoc maximum search. Exact wet-lab anchors established the performance range
of the tested hexose and process-control branches. Partial trend evidence
identified ribose as the strongest short-window pentose signal available before
validation, because it already approached the target region within the tested
3~h setting. Sequence-level priors from the BepiPred-3.0 and NetGlycate analyses
provided mechanistic context by indicating local proximity between predicted
epitope-rich regions and glycation-prone lysine positions, but this evidence was
treated as supportive and not as proof of site-specific causality. Human
constraints, reaction duration, process burden and potential overprocessing risk
were then applied as secondary feasibility filters.

On this basis, the agent team selected ribose as the final donor and locked the
pre-validation recommendation before the focused validation experiment. The locked
condition was ribose at 60~$^\circ$C for 4~h under ultrasound-assisted dry-state
glycation, with $a_w$ 0.79, a protein:sugar ratio of 1:2 (w/w), and 20~mM
phosphate buffer at pH~7.3 used in the pretreatment step. The ultrasound setting
recorded with the recommendation was 400~W for 12~min. The agent-adjudicated
central estimate was 86\% IgG-binding reduction, with a plausible range of
82--93\%. This range was treated as a heuristic evidence-weighted interval rather
than a statistically calibrated prediction interval.

Recommendation locking was used to prevent validation leakage. Once the final
ribose condition and its estimate had been recorded, the subsequent focused
validation experiment was treated as a post-recommendation test rather than an
input to the recommendation. The observed 85.6\% $\pm$ 1.7\% IgG-binding
reduction was therefore compared with the locked estimate only after the
recommendation had been fixed. The agreement between the observed value and the
agent-adjudicated range was interpreted as support for the prioritization logic,
not as evidence that the agent system had produced a calibrated quantitative
model. The additional donor-by-time confirmation panel was also treated as
post-recommendation evidence and was not used to retroactively alter the locked
decision.

\subsection*{Supplementary Methods 6. Glycation materials and experimental procedure}

The wet-lab glycation procedure was reported at a methods-reporting level rather
than as a complete bench operating protocol. The purpose of this section is to
define the experimental backbone, the variables represented in the run-sheet and
the parameters required to interpret the IgG-binding results. Routine operating
details followed previously reported protein--sugar conjugation and $\beta$-LG
dry-glycation procedures, with minor modifications for the present $\beta$-LG
system.

$\beta$-LG was dissolved in 20~mM phosphate buffer at pH~7.3 to a concentration
of 10~mg~mL$^{-1}$. Reducing sugars were added according to the run assignment
at a protein:sugar mass ratio of 1:2 (w/w). The protected run-sheet included
lactose, glucose, galactose, mannose, arabinose and ribose as donor-containing
conditions, together with no-sugar process controls. The central confirmed
materials were $\beta$-LG (L3908), ribose (R1757), glucose (G8270), galactose
(G0750) and lactose (1356676) from Sigma-Aldrich; other routine reagents were of
analytical grade unless otherwise specified. The principal reaction backbone was
dry-state glycation at $a_w$ 0.79. Reaction temperature, reaction time, donor
identity and ultrasound status followed the R01--R24 run-sheet logic described
in Supplementary Methods~4 and Supplementary Tables~3--4.

For ultrasound-assisted conditions, the dissolved $\beta$-LG sample was treated
with an ultrasonic disruptor (JY98-IIIDN, Scientz Biotechnology, Ningbo, China)
for 12~min in pulse mode, with a 2~s on/2~s off duty cycle and a nominal setting
of 400~W. The sample temperature during ultrasound treatment was maintained at
4~$^\circ$C. After pretreatment, samples were lyophilized and incubated under
the specified dry-state condition. The $a_w$ 0.79 dry-state environment was
maintained in a sealed container using saturated NaCl solution. After incubation,
reactions were stopped by reconstitution. Residual sugars and low-molecular-
weight components were removed by dialysis using a 3500-Da molecular-weight-
cutoff membrane for 48~h at 4~$^\circ$C, and the dialysed products were
lyophilized before IgG-binding analysis.

Native untreated $\beta$-LG was retained as the assay reference for IgG-binding
reduction, whereas no-sugar heated or ultrasound-treated process controls were
used to separate donor-associated glycation effects from the common processing
background. The method therefore supports comparator-aware interpretation rather
than a general claim that all reaction variables were independently optimized.

\subsection*{Supplementary Methods 7. IgG-binding assay, data processing and statistical analysis}

IgG-binding ability was measured by indirect competitive ELISA adapted from a
previous $\beta$-LG immunoreactivity method. Native $\beta$-LG was coated on
96-well microplates overnight at 4~$^\circ$C at 2~$\mu$g~mL$^{-1}$. After
washing with PBST, plates were blocked with 1\% fish-skin gelatin at
37~$^\circ$C for 1~h. For the competitive binding step, 50~$\mu$L of rabbit
anti-$\beta$-LG antiserum and 50~$\mu$L of diluted $\beta$-LG sample were added
and incubated at 37~$^\circ$C for 1~h. The diluted sample concentration used in
the competitive step was 4~$\mu$g~L$^{-1}$.

After washing, 100~$\mu$L of goat anti-rabbit IgG-HRP conjugate (LF108, Shanghai
Epizyme Biomedical Technology Co., Ltd.) was added and incubated at 37~$^\circ$C
for 1~h. The dilution ratios of rabbit antiserum and goat anti-rabbit IgG-HRP
conjugate were 1:128,000 and 1:4,000, respectively. After a further wash,
100~$\mu$L of TMB substrate was added and developed for 15~min. The reaction was
stopped with 50~$\mu$L of 2~mol~L$^{-1}$ sulfuric acid, and absorbance was
measured at 450~nm using a microplate reader (INFINITE 200 PRO, Longyue
Biotechnology Development Co., Ltd., Beijing, China). OD$_{450}$ values were
converted to IgG-binding capacity using a standard curve before reduction values
were calculated.

IgG-binding reduction was calculated uniformly relative to native untreated
$\beta$-LG using
\[
  \mathrm{Reduction} = \left[1 -
  \left(B_{\mathrm{treated}}/B_{\mathrm{native}}\right)\right] \times 100\%,
\]
where $B_{\mathrm{treated}}$ and $B_{\mathrm{native}}$ denote the standard-curve-
derived IgG-binding capacities of the treated sample and native reference,
respectively. Each condition was measured using four replicate measurements, and
each sample was analysed in four repeated ELISA wells before condition-level
values were summarized as mean $\pm$ s.d. The reported ``$\pm$'' values denote
standard deviation. No outlier exclusion was applied because no anomalous values
were identified. Statistical processing was performed using SPSS software. The
reduction endpoint is an antibody-binding immunoreactivity readout and should
not be interpreted as a clinical allergy endpoint.
\subsection*{Supplementary Methods 8. Literature-derived condition-envelope calculation}

The initial condition-envelope value shown in Fig.~4a of the main text was calculated as a conceptual upper bound on the literature-visible variable space. It was not calculated from raw JSON labels directly and was not interpreted as a number of experiments that would otherwise have been performed. Instead, raw condition labels in the structured literature evidence were first harmonized into semantically grouped levels so that synonymous labels, spelling variants, donor variants or nearby numeric settings did not inflate the apparent search space.

Six condition dimensions were used for this calculation: donor branch, reaction-time window, temperature band, pretreatment family, protein:sugar-ratio band and environment-control state. The grouped level counts were 19, 11, 7, 8, 8 and 12, respectively. The resulting conceptual upper bound was therefore
\[
19 \times 11 \times 7 \times 8 \times 8 \times 12 = 1{,}123{,}584.
\]
The grouped donor count was derived by collapsing 32 raw sugar labels into 19 donor branches, for example by merging \texttt{D-ribose} with \texttt{ribose}, \texttt{maltose} with \texttt{maltose (Mal)}, \texttt{galactose} with \texttt{galactose (Gal)}, galacto-oligosaccharide variants into one branch and dextran molecular-weight variants into a single dextran branch. Numeric reaction times and temperatures were grouped into reaction windows and processing bands, respectively, rather than counted as isolated raw values. Ratio and environment-control entries were similarly grouped into comparable bands or states, while pretreatment entries were grouped by process family. The full axis-level calculation is summarized in Supplementary Table~8.


